\section{Intended Market of Product}

The intended market of the proposed product consists of educational institutions that conduct supervised practical examinations in computer laboratories and require stronger operational control, monitoring, and post-exam accountability.

\subsection{Target Audience}
The primary target audience for the proposed system comprises higher education institutions, specifically departments and faculties offering degrees in Computer Science (CS), Software Engineering (SE), Information Technology (IT), Cybersecurity, Data Science, and other computing-intensive academic disciplines. In these fields, academic evaluation relies heavily on practical programming tasks, database design, network configurations, and software development projects conducted under controlled laboratory environments. Unlike standard theoretical assessments that can be easily managed via general-purpose Learning Management Systems (LMS) or basic multiple-choice testing tools, technical practical examinations require student interaction with complex development environments, compilers, and local operating system resources. Consequently, these departments require a specialized solution that can monitor workstation-level activities, prevent unauthorized code sharing, and secure the execution environment without hindering the technical tools required for the exam \cite{eassess1}.

Furthermore, within these academic departments, the target audience is divided into several distinct user personas, each having unique operational requirements that the system addresses. University administrators and department heads require high-level coordination features to allocate physical laboratory resources, manage lab technician roles, and maintain long-term records of exam integrity for accreditation audits. Course instructors and examiners require intuitive interfaces to configure exam sessions, upload skeleton code, define resource limitations, and compile submission archives. Invigilators and lab proctors require immediate, real-time classroom visualization tools to detect focus shifts or hardware anomalies on student machines during live testing, while students themselves require a reliable, distraction-free interface that guarantees their code submissions are securely transmitted and logged without data loss.

\subsection{Institutional Environment}
Beyond formal degree-granting university departments, the institutional environment for the proposed system extends to professional IT training academies, specialized cybersecurity bootcamps, vocational technical institutes, and official certification centers. These organizations frequently conduct high-stakes certification exams and standardized technical evaluations that demand strict compliance with security and procedural guidelines. In such professional environments, the validity of a certificate depends entirely on the integrity of the examination process. Consequently, these centers require an administrative environment that enforces strict machine allocation, verifies candidate identity, prevents external communication, and generates detailed event logs to defend against potential grading disputes or credential fraud \cite{mit_cbt}.

Additionally, the physical and network infrastructure of computer laboratories in these diverse institutional settings presents unique operational challenges that the system is engineered to navigate. Typical lab environments are characterized by shared workstation hardware, diverse software installations, and varying local area network (LAN) configurations. Managing examinations in such spaces often involves manual configuration of each workstation, creating significant overhead for IT support staff. The proposed system fits naturally into these environments by offering a lightweight, easily deployable client-server architecture that minimizes local installation complexity, operates efficiently over standard intranet networks, and provides robust local monitoring that remains resilient even during temporary network disruptions. By catering to the specific constraints of shared physical laboratories, the system provides a specialized alternative to broad, cloud-based proctoring suites.
